\begin{longtblr}[
    caption = {Параметры исследуемого образца и установки},
    label = {tbl:params},
]{
    colspec = {X[2,l] X[1,c] X[1,c]},
    hlines, vlines,
}
    Наименование параметра & Обозначение & Значение \\
    Число витков первичной обмотки & $w_1$ & \\
    Число витков вторичной обмотки & $w_2$ & \\
    Средняя длина магнитной линии, м & $l_{cp}$ & \\
    Сечение образца, м$^2$ & $S$ & \\
    Сопротивление, Ом & $R_T$ & \\
    Сопротивление интегратора, кОм & $R_И$ & \\
    Емкость интегратора, мкФ & $C_И$ & \\
    Масштаб по оси X, В/дел & $m_x$ & \\
    Масштаб по оси Y, В/дел & $m_y$ & \\
    Плотность материала, кг/м$^3$ & $d$ & \\
\end{longtblr}

\begin{longtblr}[
    caption = {Данные для построения основной кривой намагничивания ($f = 50$ Гц)},
    label = {tbl:mag_curve},
]{
    colspec = {c *{8}{X[c]}},
    hlines, vlines,
}
    $X$, дел. & 0.5 & 1.0 & 1.5 & 2.0 & 2.5 & 3.0 & 3.5 & 4.0 \\
    $Y$, дел. & & & & & & & & \\
\end{longtblr}

\begin{longtblr}[
    caption = {Данные для определения потерь энергии},
    label = {tbl:losses},
]{
    colspec = {Q[c] X[c] X[c] X[c] X[c] X[c]},
    hlines, vlines,
}
    $f$, Гц & 50 & 200 & 400 & 600 & 800 \\
    $S_п$, мм$^2$ & & & & & \\
\end{longtblr}
\newpage
\begin{longtblr}[
    caption = {Данные для определения эффективной магнитной проницаемости ($U_R = 30$ мВ)},
    label = {tbl:mu_eff},
]{
    colspec = {l *{8}{X[c]}},
    hlines, vlines,
}
    $f$, Гц & 50 & 75 & 100 & 150 & 200 & 400 & 600 & 800 \\
    $U_{BX}$, мВ & & & & & & & & \\
\end{longtblr}